\chapter{Implementation}

\section{Technology Stack}

Every technology in Table~\ref{tab:stack} was chosen for a reason that shows up in
the design: a managed runtime for the transactional backend, a shared language for
the phone and the API contract, and a toolchain that can be installed on a student
laptop. Version and usage details for the stack are taken from the official
documentation cited in the bibliography
\cite{springboot,springsecurity,postgres,expo,bleplx,esp32core,arduinojson,esp32servo,compose,jacoco,junit,mockito,jjwt}.

\begin{table}[htbp]
  \centering
  \small
  \begin{tabularx}{\textwidth}{@{}lllX@{}}
    \toprule
    \textbf{Layer} & \textbf{Technology} & \textbf{Version} & \textbf{Why} \\
    \midrule
    Backend & Spring Boot & 4.1 & Transactional REST, dependency injection and
      security out of the box. \\
    Language & Java & 21 & Records, pattern matching and a mature JPA ecosystem. \\
    Persistence & Spring Data JPA / Hibernate & Boot-managed & Repository interfaces,
      pessimistic locking, JSON column mapping. \\
    Database & PostgreSQL & 16 & Strong constraints, \code{jsonb} for pricing,
      reliable locking. \\
    Security & Spring Security + JJWT & 0.12.x & Filter chain and signed tokens
      without a custom crypto implementation. \\
    Build & Maven (wrapper) & Boot-managed & Reproducible builds in the image; the
      wrapper removes a host install. \\
    Testing & JUnit 5, Mockito, AssertJ & Boot-managed & Fast service tests with
      replaceable collaborators. \\
    Coverage & JaCoCo & 0.8.13 & Line gate scoped to the service layer. \\
    Mobile & Expo / React Native & 57 / 0.86 & One TypeScript codebase for iOS and
      Android; native module support for BLE. \\
    Mobile language & TypeScript & 6.0 & The API envelope and DTOs are typed once
      and checked at build time. \\
    BLE & \code{react-native-ble-plx} & 3.5 & Maintained GATT client with notify
      support; works in a development build. \\
    Firmware & Arduino (ESP32 core) & esp32 board package & BLE and Wi-Fi libraries
      on the same board; no RTOS work required. \\
    Firmware libs & ArduinoJson, ESP32Servo & 7.x, current & Small JSON documents,
      non-blocking servo control. \\
    Packaging & Docker, Compose v2 & 24+ & One command to build, run and test the
      backend. \\
    API testing & Postman / Newman & current & Collection-based contract checks for
      the auth controller. \\
    Version control & Git / GitHub & current & Branch-based work and history. \\
    \bottomrule
  \end{tabularx}
  \caption{Technology stack.}
  \label{tab:stack}
\end{table}

\section{Module Description}

\subsection{Identity module}

Registration hashes the password with BCrypt, ignores any client-supplied role and
writes \code{USER}. Login verifies the hash and signs a JWT whose subject is the
user id; the token carries no role and no profile, so it cannot be used to escalate
privileges and profile changes do not invalidate it. A filter converts the bearer
header into an authenticated principal whose name is that id. It parses only tokens
with three dot-separated segments, because device secrets arrive in the same header;
an invalid token is left anonymous and the authorization rules decide the response.

\textbf{Design note.} The authentication namespace is public by design ---
registration and login cannot require a token. \code{GET /auth/me} is meaningful only
with a bearer token, and it is the mobile client's launch call that turns a persisted
token back into a profile.

\subsection{Vehicle module}

Vehicle operations are owner-scoped. The controller takes the user id from the token,
never from the body, and the service resolves the record and checks ownership before
any mutation. Plate and class are the only attributes, which keeps the form short and
matches what a gate actually needs.

\subsection{Facility and administration module}

Facilities are read by \code{facilityCode}, because that is what the gate advertises.
The read model includes the pricing policy, so the application can show an hourly
estimate without a second call. Creation, pricing and status changes require the
administrator header. A new facility starts \code{INACTIVE} and must be activated
before it accepts sessions, which prevents a half-configured site from taking
parking sessions. A startup seeder inserts one active demo facility so a fresh stack
is usable immediately.

\subsection{Session module}

Session creation performs the full validation sequence in one transaction: the user
must exist, the facility must exist and be active, the vehicle must exist and belong
to the caller, and neither the user nor the vehicle may hold an active session. The
read-then-write race is closed by locking the user row, so two concurrent requests
serialise and the second sees the first session. The listing endpoint returns the
caller's own sessions, newest first, with an optional status filter and a page size
clamped to one hundred.

\subsection{Billing and payment module}

The live fee returns the frozen amount once one exists, refuses a session that has not
entered, and otherwise multiplies the facility rate by whole elapsed seconds. The
final fee is recomputed at payment and rounded up to the paisa, then written to
\code{totalFee}. Payment locks the session row, so a double tap or a client retry
sees \code{PAYMENT\_PENDING} and is refused rather than charged twice. The mock
provider implements a two-method interface and always succeeds; replacing it with a
gateway means implementing the same interface.

\subsection{Gate authorisation module}

The gate service is the only place where entry and exit are decided. It checks the
device, then the facility, then the session state, and it writes the timestamp that
starts or ends billing. Granting an entry is the only place that writes
\code{entryTimestamp}; granting an exit is the only place that writes
\code{exitGateOpenedAt}. Because both methods are transactional, a verdict and the
state it authorises cannot drift apart.

\subsection{Device and command module}

Registration generates 256 bits from a secure random source, returns the base64url
secret once and stores only its BCrypt hash. Device authentication treats an unknown
device, an unprovisioned device and a wrong secret identically, so a caller cannot
probe which ids exist. Remote commands are rows rather than sockets: an operator
queues \code{OPEN} or \code{KEEP\_CLOSED} as \code{PENDING}, the gate's next poll
takes the oldest pending row, marks it \code{DELIVERED} and returns it. The poll is
the acknowledgement, so a gate that was offline still receives the command when it
returns and never receives it twice.

\subsection{Scheduled maintenance module}

A scheduled component runs every minute and aborts \code{CREATED} sessions older than
fifteen minutes, measured against \code{createdAt}. It is a thin trigger; the query
and state change live in the session service so they can be tested without waiting
for a scheduler. Both the interval and the age threshold are configuration
properties.

\subsection{Mobile application modules}

The REST client owns the base URL, the token header and the envelope. It collapses
the backend's semantics into three outcomes every screen can handle: success, a
business refusal to display, and an authentication failure that clears the token and
returns to login. The BLE client scans by service UUID, derives the facility code
from the local name, caches one connection per facility, subscribes to the status
characteristic before writing, and resolves with the first notification that matches
the session it is waiting for. Notifications for other sessions are ignored, which
makes a concurrent remote command harmless to a waiting screen.

\subsection{Firmware modules}

The sketch is a loop around three interchangeable pieces: a transport
(\code{WiFiBackendTransport} or \code{SerialBackendTransport}), an authentication
service (\code{RemoteAuthService} or the bench-only
\code{PermissiveAuthService}), and an actuator (\code{ServoGateActuator} or the
legacy \code{LedGateActuator}). The transport is selected once at boot by sampling
GPIO26. Above those interfaces there is one gate state machine, so changing the
transport for a demonstration changes no timeout, no payload and no fail-closed
behaviour.

\reportimage[0.7]{\ImgHardwareWiring}{Gate wiring: servo signal on GPIO27, an
external 5\,V servo supply with a common ground and a decoupling capacitor, and the
transport-select jumper on GPIO26.}{fig:wiring}

\section{Key Code Excerpts}

The complete implementation is in the public repository; the excerpts below are the
four pieces that carry the most design weight.

\subsection{Session creation under a row lock}

\begin{lstlisting}[style=java,caption={The active-session check is read-then-write,
so the user row is locked for the transaction.},label={lst:create-session}]
@Transactional
public ResponseObject<UUID> createSession(SessionCreationRequestDto dto) {
    User user = userRepository.findByIdForUpdate(dto.userId()).orElse(null);
    if (user == null) return ResponseObject.failure("User does not exist.");

    ParkingFacility facility = facilityRepository
            .findByFacilityCode(dto.facilityCode()).orElse(null);
    if (facility == null)
        return ResponseObject.failure("Parking facility does not exist.");
    if (facility.getStatus() != ParkingFacilityStatus.ACTIVE)
        return ResponseObject.failure("Parking facility is not active.");

    Vehicle vehicle = vehicleRepository.findById(dto.vehicleId()).orElse(null);
    if (vehicle == null) return ResponseObject.failure("Vehicle does not exist.");
    if (!vehicle.getUser().getUserId().equals(user.getUserId()))
        return ResponseObject.failure("Vehicle does not belong to the user.");

    if (sessionRepository.existsByUser_UserIdAndSessionStatusIn(
            user.getUserId(), ACTIVE_STATUSES))
        return ResponseObject.failure("User already has an active parking session.");
    if (sessionRepository.existsByVehicle_VehicleIdAndSessionStatusIn(
            vehicle.getVehicleId(), ACTIVE_STATUSES))
        return ResponseObject.failure("Vehicle already has an active parking session.");

    ParkingSession session = SessionMapper
            .creationRequestToSession(dto, user, facility, vehicle);
    sessionRepository.save(session);
    return ResponseObject.success("Session successfully created", session.getSessionId());
}
\end{lstlisting}

\subsection{Payment freezes the fee}

\begin{lstlisting}[style=java,caption={The session row is locked, the fee is frozen
and the receipt records the backend's payment time.},label={lst:pay}]
@Transactional
public ResponseObject<PaymentReceiptDto> pay(UUID sessionId, UUID userId) {
    ParkingSession session = sessionRepository.findByIdForUpdate(sessionId)
            .orElse(null);
    if (session == null)
        return ResponseObject.failure("Parking session does not exist.");
    if (userId != null && !isOwnedBy(session, userId))
        return ResponseObject.failure("Parking session does not belong to the user.");
    if (session.getSessionStatus() != SessionStatus.PARKED)
        return ResponseObject.failure("Only a parked session can be paid for.");

    BigDecimal fee = calculateFinalFee(session);
    PaymentOutcome outcome = paymentService.charge(session.getSessionId(), fee);
    if (!outcome.successful())
        return ResponseObject.failure("Payment was not successful.");

    Instant paidAt = Instant.now();               // backend clock only
    session.setTotalFee(fee);
    session.setSessionStatus(SessionStatus.PAYMENT_PENDING);
    sessionRepository.save(session);
    // ... persist the Payment row and return the receipt
}
\end{lstlisting}

\subsection{Exit verdict with a grace re-open}

\begin{lstlisting}[style=java,caption={Exit is decided from state and payment; the
grace window is measured from the gate opening.},label={lst:exit}]
switch (session.getSessionStatus()) {
    case ABORTED -> { return keepClosed(DenyReason.STALE_SESSION); }
    case EXITED -> {
        return withinExitGraceWindow(session)
                ? open()
                : keepClosed(DenyReason.STALE_SESSION);
    }
    case PAYMENT_PENDING -> {
        if (!hasSuccessfulPayment(session))
            return keepClosed(DenyReason.PAYMENT_REQUIRED);
        Instant exitedAt = Instant.now();
        session.setSessionStatus(SessionStatus.EXITED);
        session.setExitTimestamp(exitedAt);
        session.setExitGateOpenedAt(exitedAt);
        sessionRepository.save(session);
        return open();
    }
    // CREATED / PARKED - not paid yet
    default -> { return keepClosed(DenyReason.PAYMENT_REQUIRED); }
}
\end{lstlisting}

\subsection{BLE round trip in the application}

\begin{lstlisting}[style=typescript,caption={Subscribe before write, filter by session
id, and give the whole round trip a fifteen-second budget.},label={lst:ble}]
const status = await this.withStatusNotification(device, sessionId, async () => {
  await device.writeCharacteristicWithResponseForService(
    SERVICE_UUID,
    COMMAND_CHAR_UUID,
    Buffer.from(JSON.stringify({ sessionId, intent }), "utf8").toString("base64")
  );
});

if (!status) {
  throw new Error(
    `Gate ${facilityCode} did not answer over BLE. Move closer and check power.`);
}
return { result: status.result, reason: status.reason || "NONE" };
\end{lstlisting}

\subsection{Firmware transport selection}

\begin{lstlisting}[style=cpp,caption={One pin decides the transport at boot; all gate
logic above it is shared.},label={lst:firmware-select}]
static BackendTransport *selectTransport() {
    pinMode(NETWORK_MODE_PIN, INPUT_PULLUP);
    delay(10);
    if (digitalRead(NETWORK_MODE_PIN) == LOW) return &serialTransport;
    return &wifiTransport;                       // production path
}
\end{lstlisting}
